% ============================================================
% 5. Discussion
% ============================================================

\section{Discussion}
\label{sec:discussion}

\paragraph{Key findings.}
BIC achieves its design goal: cache-bounded memory with graceful quality
degradation.  Even at extreme cache pressure (1.9\% retention,
\S\ref{sec:stress}), BIC maintains 100\% query success and 0.619
reconstruction quality---a $32\times$ improvement over baselines
that lose access to evicted agents entirely.  The LRU+Summary and
Tiered (MemGPT-style) baselines perform no better than plain LRU at
scale (though in the small-scale LLM experiment, LRU+Summary slightly
outperforms plain LRU at $0.242$ vs.\ $0.204$ semantic similarity,
suggesting summarization helps when agents produce richer content):
when children are evicted via LRU, their parents are often
already gone, so summarization fails silently.  BIC's Cantor-pairing
addressing keeps parents cached longer, enabling the summarization
chain that is the primary quality driver ($5.2\times$ improvement
per the ablation in \S\ref{sec:ablation}).  Hilbert-locality
clustering is an optimization whose benefits will materialize with
real-valued state distributions.

\paragraph{Gap between theoretical and empirical preservation ratio.}
The Information Preservation Ratio (Definition~\ref{def:preservation})
is defined as an infimum over all possible parent and child states,
providing a worst-case guarantee: no matter what states are
encountered, at least a $(1-\epsilon)$ fraction of information
survives each summarization step.  In our experiments, we report the
\emph{average} preservation ratio across all eviction events rather
than the infimum.  This is a practical choice: computing the true
infimum would require evaluating $\summarize$ over the entire
(infinite) state space, which is intractable.  The average provides
a meaningful empirical characterization of typical system behavior,
while the theoretical infimum bound provides the safety guarantee
that reconstruction quality degrades gracefully even in the worst
case.  We note that the observed average $\preserv$ values
(0.69--0.74 in our experiments) are consistent with the theoretical
framework: they represent the typical case, while the worst-case
bound ensures no catastrophic information loss for any individual
eviction event.

\paragraph{Limitations.}
We identify four limitations of the current system:

\begin{enumerate}
\item \textbf{Registry growth.}  While cached states are bounded by
  $\slots$, the agent registry grows as $\Oh(N)$.  Each record is
  lightweight ($\sim$100 bytes per agent: ID, parent ID, depth,
  status), so the registry remains small relative to state storage
  (1\,MB for 10{,}000 agents vs.\ 40\,MB for $\slots{=}1{,}000$ states).
  Nonetheless, a complete solution would bound the registry as well,
  perhaps via a secondary eviction policy or external storage.

\item \textbf{Lossy summarization.}  Hierarchical summarization is
  inherently lossy.  The preservation ratio $\preserv$ quantifies this
  loss, but the \emph{quality} of reconstructed information depends on
  the summarization function.  Our default summarizer stores raw child
  states under parent keys; a production system would use an LLM to
  produce semantic summaries.

\item \textbf{Cantor depth truncation.}  We truncate Cantor addresses
  at $d_{\max} = 8$ levels to prevent integer overflow.  With branching
  factor $k{=}3$, this supports trees of up to $3^8 = 6{,}561$ agents.
  Deeper trees require either extended-precision arithmetic or
  alternative addressing schemes.

\item \textbf{Limited LLM evaluation.}  We validate with a real LLM
  experiment (\S\ref{sec:llm}), but the scale is small (7~agents,
  Gemini-2.0-flash).  Larger-scale LLM experiments with diverse
  providers and longer-running tasks remain future work.  Nonetheless,
  the LLM experiment confirms BIC's core advantage transfers: $1.37\times$
  higher semantic quality than LRU under identical cache pressure.
\end{enumerate}

\paragraph{Hilbert locality: a non-contribution in current experiments.}
We are transparent that the ablation (Table~\ref{tab:ablation}) showed
\emph{no measurable benefit} of Hilbert-locality clustering:
BIC-Full and BIC-No-Hilbert achieve identical reconstruction quality,
semantic quality, and preservation ratio in both synthetic and LLM
settings.  We hypothesize this is because both our synthetic and
small-scale LLM experiments produce state vectors that lack sufficient
spatial structure for Hilbert curves to exploit---synthetic states are
generated with uniform random tokens, and the 7-agent LLM experiment
is too small for clustering effects to manifest.  Hilbert locality
should therefore be understood not as a current empirical contribution
but as an architectural provision for future workloads where agent
states exhibit spatial correlations (e.g., agents working on related
sub-topics that share embedding similarity).  In such settings,
Hilbert clustering would co-evict nearby agents, producing more
coherent joint summaries.  Until this hypothesis is validated on
appropriate workloads, we consider Hilbert locality a direction for
future work rather than a demonstrated advantage.

\paragraph{Broader impact.}
This work addresses a systems-level challenge in scaling agentic AI.
As LLM-based agents are deployed in increasingly complex workflows
(software engineering, scientific research, data analysis), unbounded
memory growth becomes a practical barrier to reliability and cost
efficiency.  BIC provides a principled alternative to ad-hoc memory
limits, enabling long-running multi-agent systems that operate within
predictable resource envelopes.

On the positive side, BIC makes large agent swarms more
resource-efficient by bounding cache memory, reducing the hardware
requirements for running multi-agent systems and potentially lowering
the barrier to entry for researchers and practitioners with limited
compute budgets.

However, several concerns deserve consideration.  First,
\emph{computational cost of summarization}: BIC's eviction-driven
summarization may invoke LLM calls during each eviction event (when
using LLM-powered summarization rather than raw concatenation),
adding inference cost that scales with the eviction rate.  In
high-churn scenarios, these summarization calls could dominate the
overall computational budget.  Second, \emph{environmental impact}:
by enabling larger and longer-running multi-agent systems, BIC could
contribute to increased aggregate LLM inference demand and associated
energy consumption, even as it reduces per-system memory waste.
Third, \emph{reduced human oversight}: by enabling agent swarms to
run indefinitely without crashing from resource exhaustion, BIC
removes a natural stopping point that currently forces human
intervention.  Longer-running autonomous agent systems may require
new mechanisms for periodic human review, intervention triggers, and
shutdown policies to ensure alignment with human intent.

We believe the net impact is positive---resource efficiency and
reliability are prerequisites for responsible deployment of agentic
systems---but emphasize that BIC should be deployed alongside
appropriate oversight mechanisms.


% ============================================================
\section{Conclusion}
\label{sec:conclusion}

We presented \emph{Bounded Infinity Cache} (BIC), a cache-bounded
memory management architecture for agentic swarms that provides four
formal guarantees: bounded cache consumption, zero fragmentation,
universal retrievability, and indefinite liveness.  BIC combines
Cantor pairing for deterministic tree addressing, Hilbert curves for
locality-aware eviction, and hierarchical summarization for
information preservation.  On synthetic research-decomposition tasks
with six baselines (including LRU+summarization and a Tiered
MemGPT-style baseline), BIC achieves 100\% query success and
$32\times$ higher reconstruction quality under identical cache
constraints.  A real-LLM experiment with Gemini-2.0-flash confirms
the advantage transfers: BIC achieves $1.37\times$ higher semantic
reconstruction quality than LRU.

\paragraph{Future work.}
Natural extensions include: (1)~LLM-powered semantic summarization
(replacing raw concatenation); (2)~bounding the agent registry;
(3)~distributed BIC across multiple nodes; (4)~formal verification
in Lean~4; and (5)~larger-scale LLM evaluation on production agentic
workloads.
